% Part: first-order-logic
% Chapter: syntax-and-semantics
% Section: covered-structures

\documentclass[../../../include/open-logic-section]{subfiles}

\begin{document}

\olfileid{fol}{syn}{cov}

\olsection{\printtoken{P}{structure} kang Katutup dening Term kanggo Basa Logika Tataran Kapisan}

\begin{explain}
Elinga yen sawijining term iku \emph{katutup} yen ora ngemot
!!{variable}s.
\end{explain}

\begin{defn}[!!^{value} saka term katutup]
Yen $t$ iku term katutup saka basa~$\Lang L$ lan $\Struct M$
iku !!{structure} kanggo~$\Lang L$, \emph{!!{value}}~$\Value{t}{M}$
ditetepake kaya mangkene:
\begin{enumerate}
\item Yen $t$ mung !!{constant}~$c$, mula $\Value{c}{M} = \Assign{c}{M}$.
\item Yen $t$ awujud $\Atom{f}{t_1, \ldots, t_n}$, mula
  \[
  \Value{t}{M} = \Assign{f}{M}(\Value{t_1}{M}, \ldots,
  \Value{t_n}{M}).
  \]
\end{enumerate}
\end{defn}

\begin{defn}[!!{structure} kang katutup dening term]
Sawijining !!{structure} iku \emph{katutup dening term} yen saben elemen domaine
iku !!{value} saka sawenehing term katutup.
\end{defn}

\begin{ex}
Upamakna ~$\Lang L$ iku basa kang ngemot !!{constant}s $\Obj{zero}$,
$\Obj{one}$, $\Obj{two}$, \dots, !!{predicate}~$<$ rong panggonan,
lan !!{function}s $+$ lan $\times$ kang rong panggonan. Mula
!!a{structure}~$\Struct M$ kanggo~$\Lang L$ nduweni domain
$\Domain M = \{0, 1, 2, \ldots \}$ lan pangaji
$\Assign{\Obj{zero}}{M} = 0$, $\Assign{\Obj{one}}{M} = 1$,
$\Assign{\Obj{two}}{M} = 2$, lan sabanjure. Kanggo simbol relasi
$<$ rong panggonan, himpunan $\Assign{<}{M}$ iku himpunan kabeh
pasangan $\tuple{c_1, c_2} \in \Domain{M}^2$ kanthi $c_1$
luwih cilik tinimbang~$c_2$: upamane, $\tuple{1, 3} \in
\Assign{<}{M}$ nanging $\tuple{2, 2} \notin \Assign{<}{M}$.
Kanggo !!{function} $+$ rong panggonan, tetepna $\Assign{+}{M}$
kanthi cara lumrah---upamane, $\Assign{+}{M}(2,3)$ ngasilake~$5$,
lan padha uga kanggo !!{function}~$\times$ rong panggonan.
Mula, !!{value} saka $\Obj{four}$ mung~$4$, lan !!{value}
saka $\times(\Obj{two},
+(\Obj{three},\Obj{zero}))$ (utawa ing notasi infiks,
$\Obj{two} \times (\Obj{three} + \Obj{zero})$) yaiku
kaya ing ngisor iki. Cathetan koreksi sumber OLPL-071: tampilan
Inggris nulis tandha padha karo ing pungkasan larik kapisan lan
maneh ing wiwitan larik kapindho; tandha kapisan dicopot supaya
saben langkah mung nduweni siji tandha padha karo.
\begin{multline*}
\Value{\times(\Obj{two}, +(\Obj{three},\Obj{zero}))}{M} \\
\begin{aligned}
& =\Assign{\times}{M}(\Value{\Obj{two}}{M}, \Value{+(\Obj{three}, \Obj{zero})}{M})\\
& = \Assign{\times}{M}(\Value{\Obj{two}}{M}, \Assign{+}{M}(\Value{\Obj{three}}{M},
\Value{\Obj{zero}}{M})) \\
& = \Assign{\times}{M}(\Assign{\Obj{two}}{M}, \Assign{+}{M}(\Assign{\Obj{three}}{M},
\Assign{\Obj{zero}}{M})) \\
& = \Assign{\times}{M}(2, \Assign{+}{M}(3, 0)) \\
& = \Assign{\times}{M}(2, 3) \\
& = 6
\end{aligned}
\end{multline*}
\end{ex}

\begin{prob}
Apa $\Struct N$, model standar aritmetika, katutup dening term? Jlentrehna.
\end{prob}

\end{document}
